% Structural replacement for the literal (2,3,2)-fibre words.
% This replaces the 001/010/012 rows in eq:path232-cycles and the
% paragraph which justifies their layer structure; the safe-port check is
% compressed to one exceptional-case paragraph plus seven arithmetic bounds.

\paragraph{Layer notation.}
For a boundary $t=(t_0,t_1,t_2)$ and an allowed middle colour $j\ne t_1$, put
\[
 L_j(t)=[4]\setminus\{t_0,j\},\qquad
 R_j(t)=[4]\setminus\{t_2,j\},\qquad
 H_j(t)=K_{L_j(t)}\square K_{R_j(t)}.
\]
We write $uwj$ for $(u,j,w)$, and regard $H_j(t)$ as the subgraph induced
by the vertices with middle colour $j$.  The edges between two layers are
precisely the matching edges $uwj\,uwj'$ which satisfy both endpoint lists.
The following four cycles in rook graphs will be used (the entries are
pairs $a_ib_j$):
\begin{align*}
 C_{22}(a_0,a_1;b_0,b_1)&=(a_0b_0,a_1b_0,a_1b_1,a_0b_1),\\
 C_{23}(a_0,a_1;b_0,b_1,b_2)&=(a_0b_0,a_0b_1,a_1b_1,a_1b_0,a_1b_2,a_0b_2),\\
 C_{33}(a_0,a_1,a_2;b_0,b_1,b_2)&=(a_0b_0,a_0b_1,a_1b_1,a_1b_2,\\
 &\qquad a_1b_0,a_2b_0,a_2b_1,a_2b_2,a_0b_2),\\
 C_{32}(a_0,a_1,a_2;b_0,b_1)&=(a_0b_0,a_0b_1,a_1b_1,a_2b_1,a_2b_0,a_1b_0).
\end{align*}
(The displayed $C_{32}$ is a corresponding $3	imes2$ template.)
Every consecutive pair in these displays has a common
coordinate, including the closing pair; hence these are Hamilton cycles of
$K_r\square K_s$.  If one edge is deleted, the displayed cyclic word is a
Hamilton path of the layer.

For disjoint cycles or paths, a \emph{two-edge switch} replaces
$xy,x'y'$ by $xx',yy'$ whenever the two latter pairs are matching edges
between the corresponding layers.  A chain of such switches is called a
layer splice.  This operation preserves every vertex exactly once.

\begin{lemma}[Three structural good-fibre cycles]\label{lem:structural-good-cycles}
For the normalized good boundaries $001,010,012$, the cycles $D_E,D_M,D_A$
are layer splices of the following short rook-graph templates.  Write
$C_{rs}\setminus e$ for the Hamilton path obtained by deleting the edge
$e$ from the displayed $C_{rs}$ cycle, and interpret each pair $ab$ as
$abj$ in the layer with fixed middle colour $j$.
\begin{enumerate}
\item $t=001$.  In layer $j=2$ take
$Q_2=C_{22}(1,3;0,3)\setminus(320,323)$, and in layer $j=3$ take
$Q_3=C_{22}(1,2;0,2)\setminus(230,232)$.  In layer $j=1$ use the two
paths supplied by
\[
 T_{23}(2,3;0,3,2):\quad
 [313,213,210],\ [212,312,310],
\]
where in this line $ab$ abbreviates $ab1$, and $T_{23}(a_0,a_1;b_0,b_1,b_2)$ denotes the two paths
$[a_1b_1,a_0b_1,a_0b_0]$ and
$[a_0b_2,a_1b_2,a_1b_0]$.  Join the four path ends by
$323\!\text{--}313$, $210\!\text{--}230$, $232\!\text{--}212$,
and $310\!\text{--}320$.  Their union is $D_E$.
\item $t=010$.  In layer $j=0$ take
$Q_0=C_{33}(1,3,2;1,3,2)\setminus\{(101,103),(101,102)\}$ and retain
$h=101$.  In layers $j=2,3$ take
\[
 Q_2=C_{22}(2,3;1,3)\setminus(121,123),\qquad
 Q_3=C_{22}(1,2;2,1)\setminus(131,132).
\]
Join their ends and $h$ by
$103\!\text{--}123$, $121\!\text{--}h$, $h\!\text{--}131$,
and $132\!\text{--}102$.  The union is $D_M^0$; put
$D_M^1=\rho(D_M^0)$ for the path reversal $\rho(uvw)=wvu$.
\item $t=012$.  In layer $j=2$ take
$Q_2=C_{23}(1,3;0,3,1)\setminus(121,120)$; in layer $j=0$ take
$Q_0=C_{32}(1,2,3;1,3)\setminus(101,201)$.  In layer $j=3$ retain
$h=131$ and take the path
$Q_3=[231,230,130]$ (the three non-$h$ vertices of
$C_{22}(1,2;1,0)$).  Join by
$h\!\text{--}101$, $201\!\text{--}231$,
$130\!\text{--}120$, and $121\!\text{--}h$.  The union is $D_A$.
\end{enumerate}
Every displayed cross-pair has equal endpoint colours and changes only the
middle coordinate, hence is a fibre edge.  The paths are vertex-disjoint,
and the four cross-edges close them in the cyclic order shown.  This proves
Hamiltonicity and the layer descriptions by a single inspection.  Colour
permutations and $\rho$ (with endpoint-colour interchange when needed) give
all normalized images.
\end{lemma}

\paragraph{Safe ports.}
A good boundary adjacent to a bad boundary is of type $E$ for an end change or of type $M$ for a middle change; only the $E$ transition has a nonempty unsafe set.  For
$001$, the two potentially unsafe edges are
$123\!\text{--}\!323$ and $132\!\text{--}\!232$; these are precisely the two edges declared unsafe in the definition of
$P_t$ in the main text.  All other displayed edges are safe.  The $M$ and $A$ cycles have
$P_t=\varnothing$.  Thus the safe-edge check uses only intersections of
the layer paths, rather than the full fibre words.

Up to colour permutation and reversal, every one-coordinate transition
between good boundaries is, up to the stated symmetries, one of
$aab\to aac,abb,acb,bab,cab$; $aba\to aca,cba$; or
$abc\to abd,adc$.  The two transitions not represented by the arithmetic line below are as
follows.  For $aab\to bab$, the common middle-$b$ block is the prism
$K_2\square K_3$: $E$ retains one edge in each row triangle and two matching
(vertical) edges.  $M_0$ contains both row edges and $M_1$ contains both
matching edges, so either choice gives two safe common edges.  For
$aba\to cba$, in the common $H_a$ block both cycles omit the same two edges;
after this deletion they retain respectively three and five of the seven
remaining edges, giving one common edge.  In $H_c$ they retain three and two
of four edges, giving a second common edge.  These blocks are disjoint and
both special sets are empty.

Reading the retained paths in the structural lemma gives the
following safe-edge lower bounds (the notation $x+y-m$ is the elementary
intersection bound in a common $m$-edge block):
\[
\begin{aligned}
 aab\to aac &: 3+3-4=2,&\quad
 aab\to abb &: 2+2-2=2,&\quad
 aab\to acb &: 2+2-2=2,\\
 aab\to cab &: (2+2-3)+(1+1-1)=2,&\quad
 aba\to aca &: 3+3-4=2,\\
 abc\to abd &: 4+4-6=2,&\quad
 abc\to adc &: (5+5-9)+(5+5-9)=2.
\end{aligned}
\]
Each displayed identity is read directly from the paths in the structural lemma:
for example, in $aab\to cab$ the two blocks contribute
$2+2-3=1$ and $1+1-1=1$, while in $abc\to adc$ each layer contributes
$5+5-9=1$.  The two contributions in every split identity are disjoint, so the
selected cycles have at least two common safe edges.  This proves
Lemma~\ref{lem:pathsafe} with this short capacity display and no
exhaustive list of 14-, 17-, or 16-state words.
